\scanpage{86}
la sempiternelle situation de $U=X\setminus T$ sur $S$, avec une $g_i=g_{i_0}$ section de $T$ sur $I$, et je suppose donné une section rationnelle $f$ de $U$ sur $S$, qui en un point particulier $s$ de \deleted{\unsure{car.}} codim $1$ va coïncider avec $g_{i_0}$ en tt que section de $X$
\[
g_i(s)=f(s)
\]
de sorte que l'hom induit sur le groupe d'inertie $I_{s'}\subset E_{\deleted{\unsure{S|s}}\inserted{K}}$ par le composé $E_K\longrightarrow E_{U_\eta}-E_U$ (qui envoie en fait $I_{s'}$ dans $\pi$, et $\unsure{\overline m|\overset{m}{=}}$ dans les $L_i$, on l'a vu) soit non \deleted{nulle} triviale. Considérons alors le groupe de décomposition $N_{s'}$ (cf suite exacte (7), dont on a changé les notations), qui s'identifie comme le $\pi_1$ de la ``rondelle trouée'' correspondant au hensélisé (pas strict) de la jauge $\underline O_{S,s}$ --- ou mieux, le $\pi_1$ du fibré en
\margin{85}
